%
\begin{isabellebody}%
\setisabellecontext{Week{\isacharunderscore}{\kern0pt}{\isadigit{2}}{\isacharunderscore}{\kern0pt}sol}%
%
\isadelimtheory
\isanewline
%
\endisadelimtheory
%
\isatagtheory
\isacommand{theory}\isamarkupfalse%
\ Week{\isacharunderscore}{\kern0pt}{\isadigit{2}}{\isacharunderscore}{\kern0pt}sol\isanewline
\isakeyword{imports}\ Main\isanewline
\isakeyword{begin}%
\endisatagtheory
{\isafoldtheory}%
%
\isadelimtheory
%
\endisadelimtheory
%
\begin{isamarkuptext}%
\vspace{15ex}\ExerciseSheet{7}{}%
\end{isamarkuptext}\isamarkuptrue%
%
\begin{isamarkuptext}%
Important Note: Please download Isabelle from the link provided in the slides and bring a 
      laptop to the large group tutorial. Performing the proofs in Isabelle will be an integral part
      of those tutorials.%
\end{isamarkuptext}\isamarkuptrue%
%
\begin{isamarkuptext}%
Before you start, please open a new Isabelle/HOL theory file called Week\_1.thy and add the
       following to its beginning.

\noindent theory Week\_2\\
imports Main\\
begin%
\end{isamarkuptext}\isamarkuptrue%
%
\begin{isamarkuptext}%
\Exercise{Addition Commutes}%
\end{isamarkuptext}\isamarkuptrue%
%
\begin{isamarkuptext}%
Consider the function \isa{{\isacharparenleft}{\kern0pt}{\isacharplus}{\kern0pt}{\isacharparenright}{\kern0pt}} defined in Isabelle/HOL for natural numbers.

Prove the following theorem, first pen-and-paper, and then formally in Isabelle.%
\end{isamarkuptext}\isamarkuptrue%
\isacommand{lemma}\isamarkupfalse%
\ add{\isacharunderscore}{\kern0pt}commutes{\isacharcolon}{\kern0pt}\ {\isachardoublequoteopen}{\isacharparenleft}{\kern0pt}n{\isacharcolon}{\kern0pt}{\isacharcolon}{\kern0pt}nat{\isacharparenright}{\kern0pt}\ {\isacharplus}{\kern0pt}\ m\ {\isacharequal}{\kern0pt}\ m\ {\isacharplus}{\kern0pt}\ n{\isachardoublequoteclose}%
\begin{isamarkuptext}%
\paragraph{Hints:} 

 1. Perform the proof by inuction on \isa{n}, using substitutions \isa{subst}, and forward reasoning using
   the method \isa{rule}.

 2. You will need to identify multiple lemmas about how \isa{add} works. Do not prove those lemmas. Use 
    \isa{sorry} as the proof. This is the top-down approach. 

 3. There is a proof which uses the following lemmas:

   \isa{add{\isacharunderscore}{\kern0pt}{\isadigit{0}}{\isacharunderscore}{\kern0pt}right{\isacharcomma}{\kern0pt}\ add{\isacharunderscore}{\kern0pt}{\isadigit{0}}{\isacharunderscore}{\kern0pt}left{\isacharcomma}{\kern0pt}\ add{\isacharunderscore}{\kern0pt}Suc{\isacharunderscore}{\kern0pt}right{\isacharcomma}{\kern0pt}\ add{\isacharunderscore}{\kern0pt}Suc}%
\end{isamarkuptext}\isamarkuptrue%
%
\isadelimproof
%
\endisadelimproof
%
\isatagproof
\isacommand{proof}\isamarkupfalse%
{\isacharparenleft}{\kern0pt}induction\ n{\isacharparenright}{\kern0pt}\isanewline
\ \ \isacommand{case}\isamarkupfalse%
\ {\isadigit{0}}\isanewline
\ \ \isacommand{have}\isamarkupfalse%
\ {\isadigit{1}}{\isacharcolon}{\kern0pt}\ {\isachardoublequoteopen}{\isadigit{0}}\ {\isacharplus}{\kern0pt}\ m\ {\isacharequal}{\kern0pt}\ m{\isachardoublequoteclose}\isanewline
\ \ \ \ \isacommand{using}\isamarkupfalse%
\ Groups{\isachardot}{\kern0pt}comm{\isacharunderscore}{\kern0pt}monoid{\isacharunderscore}{\kern0pt}add{\isacharunderscore}{\kern0pt}class{\isachardot}{\kern0pt}add{\isacharunderscore}{\kern0pt}{\isadigit{0}}{\isacharbrackleft}{\kern0pt}\isakeyword{where}\ {\isacharquery}{\kern0pt}a\ {\isacharequal}{\kern0pt}\ m{\isacharbrackright}{\kern0pt}\isanewline
\ \ \ \ \isacommand{{\isachardot}{\kern0pt}}\isamarkupfalse%
\isanewline
\isanewline
\ \ \isacommand{have}\isamarkupfalse%
\ {\isadigit{2}}{\isacharcolon}{\kern0pt}\ {\isachardoublequoteopen}m\ {\isacharplus}{\kern0pt}\ {\isadigit{0}}\ {\isacharequal}{\kern0pt}\ m{\isachardoublequoteclose}\isanewline
\ \ \ \ \isacommand{using}\isamarkupfalse%
\ Groups{\isachardot}{\kern0pt}monoid{\isacharunderscore}{\kern0pt}add{\isacharunderscore}{\kern0pt}class{\isachardot}{\kern0pt}add{\isacharunderscore}{\kern0pt}{\isadigit{0}}{\isacharunderscore}{\kern0pt}right\isanewline
\ \ \ \ \isacommand{{\isachardot}{\kern0pt}}\isamarkupfalse%
\isanewline
\isanewline
\ \ \isacommand{show}\isamarkupfalse%
\ {\isacharquery}{\kern0pt}case\isanewline
\ \ \ \ \isacommand{apply}\isamarkupfalse%
{\isacharparenleft}{\kern0pt}subst\ {\isadigit{1}}{\isacharparenright}{\kern0pt}\isanewline
\ \ \ \ \isacommand{apply}\isamarkupfalse%
{\isacharparenleft}{\kern0pt}subst\ {\isadigit{2}}{\isacharparenright}{\kern0pt}\isanewline
\ \ \ \ \isacommand{{\isachardot}{\kern0pt}{\isachardot}{\kern0pt}}\isamarkupfalse%
\isanewline
\isanewline
\isacommand{next}\isamarkupfalse%
\isanewline
\ \ \isacommand{case}\isamarkupfalse%
\ {\isacharparenleft}{\kern0pt}Suc\ n{\isacharparenright}{\kern0pt}\isanewline
\ \ \isacommand{have}\isamarkupfalse%
\ {\isadigit{1}}{\isacharcolon}{\kern0pt}\ {\isachardoublequoteopen}{\isacharparenleft}{\kern0pt}Suc\ n{\isacharparenright}{\kern0pt}\ {\isacharplus}{\kern0pt}\ m\ {\isacharequal}{\kern0pt}\ Suc\ {\isacharparenleft}{\kern0pt}n\ {\isacharplus}{\kern0pt}\ m{\isacharparenright}{\kern0pt}{\isachardoublequoteclose}\isanewline
\ \ \ \ \isacommand{using}\isamarkupfalse%
\ Nat{\isachardot}{\kern0pt}plus{\isacharunderscore}{\kern0pt}nat{\isachardot}{\kern0pt}add{\isacharunderscore}{\kern0pt}Suc\isanewline
\ \ \ \ \isacommand{{\isachardot}{\kern0pt}}\isamarkupfalse%
\isanewline
\isanewline
\ \ \isacommand{have}\isamarkupfalse%
\ {\isadigit{2}}{\isacharcolon}{\kern0pt}\ {\isachardoublequoteopen}m\ {\isacharplus}{\kern0pt}\ Suc\ n\ {\isacharequal}{\kern0pt}\ Suc\ {\isacharparenleft}{\kern0pt}m\ {\isacharplus}{\kern0pt}\ n{\isacharparenright}{\kern0pt}{\isachardoublequoteclose}\isanewline
\ \ \ \ \isacommand{using}\isamarkupfalse%
\ Nat{\isachardot}{\kern0pt}add{\isacharunderscore}{\kern0pt}Suc{\isacharunderscore}{\kern0pt}right\isanewline
\ \ \ \ \isacommand{{\isachardot}{\kern0pt}}\isamarkupfalse%
\isanewline
\isanewline
\ \ \isacommand{show}\isamarkupfalse%
\ {\isacharquery}{\kern0pt}case\isanewline
\ \ \ \ \isacommand{apply}\isamarkupfalse%
{\isacharparenleft}{\kern0pt}subst\ {\isadigit{1}}{\isacharparenright}{\kern0pt}\isanewline
\ \ \ \ \isacommand{apply}\isamarkupfalse%
{\isacharparenleft}{\kern0pt}subst\ {\isadigit{2}}{\isacharparenright}{\kern0pt}\isanewline
\ \ \ \ \isacommand{apply}\isamarkupfalse%
{\isacharparenleft}{\kern0pt}subst\ Suc{\isachardot}{\kern0pt}IH{\isacharparenright}{\kern0pt}\isanewline
\ \ \ \ \isacommand{{\isachardot}{\kern0pt}{\isachardot}{\kern0pt}}\isamarkupfalse%
\isanewline
\isanewline
\isacommand{qed}\isamarkupfalse%
%
\endisatagproof
{\isafoldproof}%
%
\isadelimproof
%
\endisadelimproof
%
\begin{isamarkuptext}%
\Exercise{Multiplication is Monotone}%
\end{isamarkuptext}\isamarkuptrue%
%
\begin{isamarkuptext}%
Consider the multiplication function \isa{{\isacharparenleft}{\kern0pt}{\isacharasterisk}{\kern0pt}{\isacharparenright}{\kern0pt}} defined in Isabelle. Prove the following
      property of it, i.e.\ that it is monotonically increasing in both arguments. You should prove
      it both, pen-and-paper and in Isabelle.%
\end{isamarkuptext}\isamarkuptrue%
\isacommand{lemma}\isamarkupfalse%
\ mult{\isacharunderscore}{\kern0pt}mono{\isacharcolon}{\kern0pt}\isanewline
\ \ \isakeyword{fixes}\ a\ b\ c\ d{\isacharcolon}{\kern0pt}{\isacharcolon}{\kern0pt}\ nat\ %
\isamarkupcmt{Note the alternative way of fixing the types of the constants%
}\isanewline
\ \ \isakeyword{assumes}\ {\isachardoublequoteopen}a\ {\isasymle}\ b{\isachardoublequoteclose}\ {\isachardoublequoteopen}c\ {\isasymle}\ d{\isachardoublequoteclose}\isanewline
\ \ \isakeyword{shows}\ {\isachardoublequoteopen}a\ {\isacharasterisk}{\kern0pt}\ c\ {\isasymle}\ b\ {\isacharasterisk}{\kern0pt}\ d{\isachardoublequoteclose}%
\begin{isamarkuptext}%
Your pen-and-paper proof should indicate

 1. what lemmas are you assuming,

 2. on what variable are you preforming induction,

 3. what are assumptions and the proof goals in the base case and the step case, and what is the
    induction hypothesis, and

 4. how are you instantiating the induction hypothesis, i.e. how do you show that its assumptions
    hold and which quantified variables are instantiated with which constants.

\paragraph{Hints:}
 1. The proof should be by induction

 2. The following lemmas suffice to prove the goal with only substitution, forward reasoning, and
     induction:

     \isa{refl{\isacharcomma}{\kern0pt}\ diff{\isacharunderscore}{\kern0pt}le{\isacharunderscore}{\kern0pt}self{\isacharcomma}{\kern0pt}\ mult{\isacharunderscore}{\kern0pt}le{\isacharunderscore}{\kern0pt}mono{\isadigit{2}}{\isacharcomma}{\kern0pt}\ mult{\isacharunderscore}{\kern0pt}zero{\isacharunderscore}{\kern0pt}right{\isacharcomma}{\kern0pt}\ le{\isadigit{0}}{\isacharcomma}{\kern0pt}\ mult{\isacharunderscore}{\kern0pt}Suc{\isacharunderscore}{\kern0pt}right{\isacharcomma}{\kern0pt}\ add{\isacharunderscore}{\kern0pt}mono{\isacharcomma}{\kern0pt}\ Suc{\isacharunderscore}{\kern0pt}eq{\isacharunderscore}{\kern0pt}plus{\isadigit{1}}{\isacharcomma}{\kern0pt}\ le{\isacharunderscore}{\kern0pt}diff{\isacharunderscore}{\kern0pt}conv}

 3. When you have a term like \isa{x\ {\isacharminus}{\kern0pt}\ y}, where \isa{x} is a \isa{nat}, in the goal or the assumptions, you 
     should perform proof by case analysis on whether \isa{y\ {\isasymle}\ x}. This is because of the way \isa{{\isacharminus}{\kern0pt}} is
     defined for natural numbers, where \isa{x\ {\isacharminus}{\kern0pt}\ y\ {\isacharequal}{\kern0pt}\ {\isadigit{0}}}, for any \isa{x\ {\isasymle}\ y}.%
\end{isamarkuptext}\isamarkuptrue%
%
\isadelimproof
\ \ %
\endisadelimproof
%
\isatagproof
\isacommand{using}\isamarkupfalse%
\ assms\isanewline
\isacommand{proof}\isamarkupfalse%
{\isacharparenleft}{\kern0pt}induction\ d{\isacharparenright}{\kern0pt}\isanewline
\ \ \isacommand{case}\isamarkupfalse%
\ assms{\isacharcolon}{\kern0pt}\ {\isadigit{0}}\isanewline
\ \ \isacommand{have}\isamarkupfalse%
\ {\isadigit{1}}{\isacharcolon}{\kern0pt}\ {\isachardoublequoteopen}c\ {\isacharequal}{\kern0pt}\ {\isadigit{0}}{\isachardoublequoteclose}\isanewline
\ \ \ \ \isacommand{using}\isamarkupfalse%
\ \ Nat{\isachardot}{\kern0pt}bot{\isacharunderscore}{\kern0pt}nat{\isacharunderscore}{\kern0pt}{\isadigit{0}}{\isachardot}{\kern0pt}extremum{\isacharunderscore}{\kern0pt}uniqueI{\isacharbrackleft}{\kern0pt}OF\ assms{\isacharparenleft}{\kern0pt}{\isadigit{2}}{\isacharparenright}{\kern0pt}{\isacharbrackright}{\kern0pt}\isanewline
\ \ \ \ \isacommand{{\isachardot}{\kern0pt}}\isamarkupfalse%
\isanewline
\ \ \ \ \isacommand{find{\isacharunderscore}{\kern0pt}theorems}\isamarkupfalse%
\ {\isachardoublequoteopen}{\isacharquery}{\kern0pt}x\ {\isasymle}\ {\isacharquery}{\kern0pt}x{\isachardoublequoteclose}\isanewline
\isanewline
\ \ \isacommand{show}\isamarkupfalse%
\ {\isacharquery}{\kern0pt}case\isanewline
\ \ \ \ \isacommand{apply}\isamarkupfalse%
{\isacharparenleft}{\kern0pt}subst\ {\isadigit{1}}{\isacharparenright}{\kern0pt}\ \isanewline
\ \ \ \ \isacommand{apply}\isamarkupfalse%
{\isacharparenleft}{\kern0pt}subst\ mult{\isacharunderscore}{\kern0pt}zero{\isacharunderscore}{\kern0pt}right{\isacharparenright}{\kern0pt}\isanewline
\ \ \ \ \isacommand{apply}\isamarkupfalse%
{\isacharparenleft}{\kern0pt}subst\ mult{\isacharunderscore}{\kern0pt}zero{\isacharunderscore}{\kern0pt}right{\isacharparenright}{\kern0pt}\isanewline
\ \ \ \ \isacommand{using}\isamarkupfalse%
\ preorder{\isacharunderscore}{\kern0pt}class{\isachardot}{\kern0pt}order{\isachardot}{\kern0pt}refl\isanewline
\ \ \ \ \isacommand{{\isachardot}{\kern0pt}}\isamarkupfalse%
\isanewline
\isanewline
\isacommand{next}\isamarkupfalse%
\isanewline
\ \ \isacommand{case}\isamarkupfalse%
\ Suc{\isacharcolon}{\kern0pt}\ {\isacharparenleft}{\kern0pt}Suc\ d{\isacharparenright}{\kern0pt}\isanewline
\isanewline
\ \ \isacommand{show}\isamarkupfalse%
\ {\isacharquery}{\kern0pt}case\isanewline
\ \ \isacommand{proof}\isamarkupfalse%
{\isacharparenleft}{\kern0pt}cases\ {\isachardoublequoteopen}c\ {\isacharequal}{\kern0pt}\ Suc\ d{\isachardoublequoteclose}{\isacharparenright}{\kern0pt}\isanewline
\ \ \ \ \isacommand{case}\isamarkupfalse%
\ c{\isacharunderscore}{\kern0pt}eq{\isacharunderscore}{\kern0pt}sucd{\isacharcolon}{\kern0pt}\ True\isanewline
\ \ \ \ \isacommand{show}\isamarkupfalse%
\ {\isacharquery}{\kern0pt}thesis\isanewline
\ \ \ \ \ \ \isacommand{apply}\isamarkupfalse%
{\isacharparenleft}{\kern0pt}subst\ c{\isacharunderscore}{\kern0pt}eq{\isacharunderscore}{\kern0pt}sucd{\isacharparenright}{\kern0pt}\isanewline
\ \ \ \ \ \ \isacommand{using}\isamarkupfalse%
\ mult{\isacharunderscore}{\kern0pt}le{\isacharunderscore}{\kern0pt}mono{\isadigit{1}}{\isacharbrackleft}{\kern0pt}OF\ assms{\isacharparenleft}{\kern0pt}{\isadigit{1}}{\isacharparenright}{\kern0pt}{\isacharbrackright}{\kern0pt}\isanewline
\ \ \ \ \ \ \isacommand{{\isachardot}{\kern0pt}}\isamarkupfalse%
\isanewline
\isanewline
\ \ \isacommand{next}\isamarkupfalse%
\isanewline
\ \ \ \ \isacommand{case}\isamarkupfalse%
\ c{\isacharunderscore}{\kern0pt}lt{\isacharunderscore}{\kern0pt}sucd{\isacharcolon}{\kern0pt}\ False\isanewline
\isanewline
\ \ \ \ \isacommand{have}\isamarkupfalse%
\ {\isadigit{1}}{\isacharcolon}{\kern0pt}\ {\isachardoublequoteopen}c\ {\isachargreater}{\kern0pt}\ Suc\ d\ {\isasymor}\ c\ {\isacharless}{\kern0pt}\ Suc\ d{\isachardoublequoteclose}\isanewline
\ \ \ \ \ \ \isacommand{apply}\isamarkupfalse%
{\isacharparenleft}{\kern0pt}subst\ neq{\isacharunderscore}{\kern0pt}iff{\isacharbrackleft}{\kern0pt}symmetric{\isacharbrackright}{\kern0pt}{\isacharparenright}{\kern0pt}\isanewline
\ \ \ \ \ \ \isacommand{using}\isamarkupfalse%
\ c{\isacharunderscore}{\kern0pt}lt{\isacharunderscore}{\kern0pt}sucd{\isacharbrackleft}{\kern0pt}symmetric{\isacharbrackright}{\kern0pt}\isanewline
\ \ \ \ \ \ \isacommand{{\isachardot}{\kern0pt}}\isamarkupfalse%
\isanewline
\ \ \ \ \isacommand{have}\isamarkupfalse%
\ {\isadigit{2}}{\isacharcolon}{\kern0pt}\ {\isachardoublequoteopen}{\isasymnot}\ c\ {\isachargreater}{\kern0pt}\ Suc\ d{\isachardoublequoteclose}\isanewline
\ \ \ \ \ \ \isacommand{using}\isamarkupfalse%
\ leD{\isacharbrackleft}{\kern0pt}OF\ Suc{\isacharparenleft}{\kern0pt}{\isadigit{3}}{\isacharparenright}{\kern0pt}{\isacharbrackright}{\kern0pt}\isanewline
\ \ \ \ \ \ \isacommand{{\isachardot}{\kern0pt}}\isamarkupfalse%
\isanewline
\ \ \ \ \isacommand{have}\isamarkupfalse%
\ {\isadigit{3}}{\isacharcolon}{\kern0pt}\ {\isachardoublequoteopen}c\ {\isacharless}{\kern0pt}\ Suc\ d{\isachardoublequoteclose}\isanewline
\ \ \ \ \ \ \isacommand{using}\isamarkupfalse%
\ Meson{\isachardot}{\kern0pt}make{\isacharunderscore}{\kern0pt}pos{\isacharunderscore}{\kern0pt}rule{\isacharprime}{\kern0pt}{\isacharbrackleft}{\kern0pt}OF\ {\isadigit{1}}\ {\isadigit{2}}{\isacharbrackright}{\kern0pt}\isanewline
\ \ \ \ \ \ \isacommand{{\isachardot}{\kern0pt}}\isamarkupfalse%
\isanewline
\isanewline
\ \ \ \ \isacommand{have}\isamarkupfalse%
\ {\isadigit{4}}{\isacharcolon}{\kern0pt}\ {\isachardoublequoteopen}c\ {\isasymle}\ d{\isachardoublequoteclose}\isanewline
\ \ \ \ \ \ \isacommand{apply}\isamarkupfalse%
{\isacharparenleft}{\kern0pt}subst\ less{\isacharunderscore}{\kern0pt}Suc{\isacharunderscore}{\kern0pt}eq{\isacharunderscore}{\kern0pt}le{\isacharbrackleft}{\kern0pt}symmetric{\isacharbrackright}{\kern0pt}{\isacharparenright}{\kern0pt}\isanewline
\ \ \ \ \ \ \isacommand{using}\isamarkupfalse%
\ {\isadigit{3}}\isanewline
\ \ \ \ \ \ \isacommand{{\isachardot}{\kern0pt}}\isamarkupfalse%
\isanewline
\ \ \ \ \ \ \ \ \ \ \ \ \ \ \ \ \ \isanewline
\isanewline
\ \ \ \ \isacommand{have}\isamarkupfalse%
\ {\isadigit{5}}{\isacharcolon}{\kern0pt}\ {\isachardoublequoteopen}a\ {\isacharasterisk}{\kern0pt}\ c\ {\isasymle}\ b\ {\isacharasterisk}{\kern0pt}\ d{\isachardoublequoteclose}\isanewline
\ \ \ \ \ \ \isacommand{using}\isamarkupfalse%
\ Suc{\isachardot}{\kern0pt}IH{\isacharbrackleft}{\kern0pt}OF\ assms{\isacharparenleft}{\kern0pt}{\isadigit{1}}{\isacharparenright}{\kern0pt}\ {\isadigit{4}}{\isacharbrackright}{\kern0pt}\isanewline
\ \ \ \ \ \ \isacommand{{\isachardot}{\kern0pt}}\isamarkupfalse%
\isanewline
\isanewline
\ \ \ \ \isacommand{have}\isamarkupfalse%
\ {\isadigit{6}}{\isacharcolon}{\kern0pt}\ {\isachardoublequoteopen}b\ {\isacharasterisk}{\kern0pt}\ d\ {\isasymle}\ b\ {\isacharasterisk}{\kern0pt}\ {\isacharparenleft}{\kern0pt}Suc\ d{\isacharparenright}{\kern0pt}{\isachardoublequoteclose}\isanewline
\ \ \ \ \ \ \isacommand{apply}\isamarkupfalse%
{\isacharparenleft}{\kern0pt}rule\ mult{\isacharunderscore}{\kern0pt}le{\isacharunderscore}{\kern0pt}mono{\isadigit{2}}{\isacharparenright}{\kern0pt}\isanewline
\ \ \ \ \ \ \isacommand{apply}\isamarkupfalse%
{\isacharparenleft}{\kern0pt}rule\ order{\isachardot}{\kern0pt}strict{\isacharunderscore}{\kern0pt}implies{\isacharunderscore}{\kern0pt}order{\isacharparenright}{\kern0pt}\isanewline
\ \ \ \ \ \ \isacommand{using}\isamarkupfalse%
\ Nat{\isachardot}{\kern0pt}lessI\isanewline
\ \ \ \ \ \ \isacommand{{\isachardot}{\kern0pt}}\isamarkupfalse%
\isanewline
\ \ \ \ \isacommand{show}\isamarkupfalse%
\ {\isacharquery}{\kern0pt}thesis\isanewline
\ \ \ \ \ \ \isacommand{using}\isamarkupfalse%
\ le{\isacharunderscore}{\kern0pt}trans{\isacharbrackleft}{\kern0pt}OF\ {\isadigit{5}}\ {\isadigit{6}}{\isacharbrackright}{\kern0pt}\isanewline
\ \ \ \ \ \ \isacommand{{\isachardot}{\kern0pt}}\isamarkupfalse%
\isanewline
\isanewline
\ \ \isacommand{qed}\isamarkupfalse%
\isanewline
\isanewline
\isanewline
\isacommand{qed}\isamarkupfalse%
%
\endisatagproof
{\isafoldproof}%
%
\isadelimproof
%
\endisadelimproof
\isanewline
\isanewline
\isanewline
\isanewline
%
\isadelimtheory
\isanewline
%
\endisadelimtheory
%
\isatagtheory
\isacommand{end}\isamarkupfalse%
%
\endisatagtheory
{\isafoldtheory}%
%
\isadelimtheory
%
\endisadelimtheory
%
\end{isabellebody}%
\endinput
%:%file=Week_2_sol.tex%:%
%:%6=1%:%
%:%11=2%:%
%:%12=2%:%
%:%13=3%:%
%:%14=4%:%
%:%23=6%:%
%:%27=8%:%
%:%28=9%:%
%:%29=10%:%
%:%33=13%:%
%:%34=14%:%
%:%35=15%:%
%:%36=16%:%
%:%37=17%:%
%:%38=18%:%
%:%42=21%:%
%:%46=25%:%
%:%47=26%:%
%:%48=27%:%
%:%50=44%:%
%:%51=44%:%
%:%53=46%:%
%:%54=47%:%
%:%55=48%:%
%:%56=49%:%
%:%57=50%:%
%:%58=51%:%
%:%59=52%:%
%:%60=53%:%
%:%61=54%:%
%:%62=55%:%
%:%63=56%:%
%:%71=59%:%
%:%72=59%:%
%:%73=60%:%
%:%74=60%:%
%:%75=61%:%
%:%76=61%:%
%:%77=62%:%
%:%78=62%:%
%:%79=63%:%
%:%80=63%:%
%:%81=64%:%
%:%82=65%:%
%:%83=65%:%
%:%84=66%:%
%:%85=66%:%
%:%86=67%:%
%:%87=67%:%
%:%88=68%:%
%:%89=69%:%
%:%90=69%:%
%:%91=70%:%
%:%92=70%:%
%:%93=71%:%
%:%94=71%:%
%:%95=72%:%
%:%96=72%:%
%:%97=73%:%
%:%98=74%:%
%:%99=74%:%
%:%100=75%:%
%:%101=75%:%
%:%102=76%:%
%:%103=76%:%
%:%104=77%:%
%:%105=77%:%
%:%106=78%:%
%:%107=78%:%
%:%108=79%:%
%:%109=80%:%
%:%110=80%:%
%:%111=81%:%
%:%112=81%:%
%:%113=82%:%
%:%114=82%:%
%:%115=83%:%
%:%116=84%:%
%:%117=84%:%
%:%118=85%:%
%:%119=85%:%
%:%120=86%:%
%:%121=86%:%
%:%122=87%:%
%:%123=87%:%
%:%124=88%:%
%:%125=88%:%
%:%126=89%:%
%:%127=90%:%
%:%137=96%:%
%:%141=98%:%
%:%142=99%:%
%:%143=100%:%
%:%145=123%:%
%:%146=123%:%
%:%147=124%:%
%:%148=124%:%
%:%149=124%:%
%:%150=125%:%
%:%151=126%:%
%:%153=129%:%
%:%154=130%:%
%:%155=131%:%
%:%156=132%:%
%:%157=133%:%
%:%158=134%:%
%:%159=135%:%
%:%160=136%:%
%:%161=137%:%
%:%162=138%:%
%:%163=139%:%
%:%164=140%:%
%:%165=141%:%
%:%166=142%:%
%:%167=143%:%
%:%168=144%:%
%:%169=145%:%
%:%170=146%:%
%:%171=147%:%
%:%171=148%:%
%:%172=149%:%
%:%173=150%:%
%:%174=151%:%
%:%175=152%:%
%:%179=158%:%
%:%183=158%:%
%:%184=158%:%
%:%185=159%:%
%:%186=159%:%
%:%187=160%:%
%:%188=160%:%
%:%189=161%:%
%:%190=161%:%
%:%191=162%:%
%:%192=162%:%
%:%193=163%:%
%:%194=163%:%
%:%195=164%:%
%:%196=164%:%
%:%197=165%:%
%:%198=166%:%
%:%199=166%:%
%:%200=167%:%
%:%201=167%:%
%:%202=168%:%
%:%203=168%:%
%:%204=169%:%
%:%205=169%:%
%:%206=170%:%
%:%207=170%:%
%:%208=171%:%
%:%209=171%:%
%:%210=172%:%
%:%211=173%:%
%:%212=173%:%
%:%213=174%:%
%:%214=174%:%
%:%215=175%:%
%:%216=176%:%
%:%217=176%:%
%:%218=177%:%
%:%219=177%:%
%:%220=178%:%
%:%221=178%:%
%:%222=179%:%
%:%223=179%:%
%:%224=180%:%
%:%225=180%:%
%:%226=181%:%
%:%227=181%:%
%:%228=182%:%
%:%229=182%:%
%:%230=183%:%
%:%231=184%:%
%:%232=184%:%
%:%233=185%:%
%:%234=185%:%
%:%235=186%:%
%:%236=187%:%
%:%237=187%:%
%:%238=188%:%
%:%239=188%:%
%:%240=189%:%
%:%241=189%:%
%:%242=190%:%
%:%243=190%:%
%:%244=191%:%
%:%245=191%:%
%:%246=192%:%
%:%247=192%:%
%:%248=193%:%
%:%249=193%:%
%:%250=194%:%
%:%251=194%:%
%:%252=195%:%
%:%253=195%:%
%:%254=196%:%
%:%255=196%:%
%:%256=197%:%
%:%257=198%:%
%:%258=198%:%
%:%259=199%:%
%:%260=199%:%
%:%261=200%:%
%:%262=200%:%
%:%263=201%:%
%:%264=201%:%
%:%265=202%:%
%:%266=203%:%
%:%267=204%:%
%:%268=204%:%
%:%269=205%:%
%:%270=205%:%
%:%271=206%:%
%:%272=206%:%
%:%273=207%:%
%:%274=208%:%
%:%275=208%:%
%:%276=209%:%
%:%277=209%:%
%:%278=210%:%
%:%279=210%:%
%:%280=211%:%
%:%281=211%:%
%:%282=212%:%
%:%283=212%:%
%:%284=213%:%
%:%285=213%:%
%:%286=214%:%
%:%287=214%:%
%:%288=215%:%
%:%289=215%:%
%:%290=216%:%
%:%291=217%:%
%:%292=217%:%
%:%293=218%:%
%:%294=219%:%
%:%295=220%:%
%:%303=220%:%
%:%304=221%:%
%:%305=222%:%
%:%306=223%:%
%:%309=224%:%
%:%314=225%:%
